\documentclass[10pt]{article}
\usepackage[margin=20mm]{geometry}
\usepackage{amsmath,amssymb,mathtools,fontspec,unicode-math,xurl,hyperref}
\IfFontExistsTF{DejaVu Serif}{\setmainfont{DejaVu Serif}}{\setmainfont{Latin Modern Roman}}
\setmathfont{latinmodern-math.otf}
\setlength{\parindent}{0pt}\setlength{\parskip}{6pt}\sloppy
\title{Verification of the Apollonian dimension certificate}
\author{}
\date{25 October 2026}
\begin{document}
\maketitle
Abstract: This initial review examines the publisher version of Vytnova and Wormell’s Hausdorff dimension of the Apollonian gasket. A complete rational counterexample disproves the parameter enclosure in Theorem 4.2 under the specified Hardy-space interpretation of its operator assumptions. The same unsafe choice of denominator occurs in the pinned public rigorous implementation. A second problem is the asserted density of analytic functions in the usual Hölder norm in Proposition 2.3; a replacement argument restores that spectral identification. Both concerns affect the proof of the central certificate. The review does not establish that the published numerical value is false. The parameter-enclosure repair is proved here; applying it to the authors’ actual high-precision residual remains outstanding. Focused independent acceptance retains both findings; human adjudication remains pending.\par
\section*{Identity, evidence and inspected scope}
Paper. Polina L. Vytnova and Caroline L. Wormell, Hausdorff dimension of the Apollonian gasket, Inventiones mathematicae 239 (2025), 909–946, DOI https://doi.org/10.1007/s00222-024-01311-y. The reviewed document is the 38-page publisher version of record, published online on 22 January 2025 [VW]. Paper ID: M25-0051. The page convention is one-based PDF page = printed page minus 908.\par
The entire extracted main text, including references, was read. Detailed checks concentrated on the dimension-to-spectrum bridge (Proposition 2.3), the two-dimensional Chebyshev estimate and its role in error control (Proposition 3.2 and Theorem 3.3), the residual-to-parameter enclosure (Theorems 4.1–4.2), and the dependency of the final certificate on these steps (Propositions 6.4–6.5 and Section 7). The repository was pinned to commit `4f73a6a1807ea7cb28313c379ffeeaf1410ff1d6` of 13 February 2025 [code]; its identity as the exact revision used for the published run is unverified.\par
The acquired arXiv TeX archive, version 2406.04922v1 of 7 June 2024, is a separate source version; it was not substituted for the journal text. The optional `APriori.ipynb` was not acquired. Its absence is a coverage limitation and is not a mathematical finding. No author code, notebook, JLD input or 448-bit computation was executed. The review does not certify every numerical constant, interval-library operation, or every general-dimensional endpoint in Section 3. The accompanying coverage ledger distinguishes reading from independently verified reasoning.\par
\section*{Problem 1: the parameter enclosure fails for negative residuals}
Judgment: Major; demonstrated invalid theorem and certification step. The issue is in Theorem 4.2, printed pp. 928–929 (PDF pp. 20–21). Its significance is explicit in the introduction, p. 911, and the proof of Theorem 1.1, p. 943.\par
\subsection*{Claim and suspect step}
Equation reference: derivative, printed. For a positive parameterized transfer operator, the theorem assumes \[
-D^-\sup\psi\ \le\ \frac{d}{ds}(\mathcal L_s\psi)(x)\
 \le\ -D^+\inf\psi <0,
\qquad D^\pm>0.\] The source states (29) for every bounded \(\psi>0\). Here \(\mathcal L_s\) is the operator (13) on the Hardy space (11), whose restrictions to the compact real cube are continuous; every strictly positive function in that domain has positive infimum. The counterexample below satisfies (29) for every such function, and more generally for every bounded function bounded away from zero. If bounded \(\psi>0\) were instead read as all pointwise-positive bounded functions, its terminal strict inequality would be impossible for functions with zero infimum. The enclosure failure is established under the coherent Hardy-space (or uniformly positive bounded-function) interpretation, without relying on that wording ambiguity. For a trial function \(0<\phi^-\le\phi\le\phi^+\), it assumes \[\varepsilon^-<(\mathcal P_K\mathcal L_{s_0}\phi-\lambda\phi)(x)<\varepsilon^+,
 \qquad
 \mathrm{err}=\|\mathcal L_{s_0}\phi-\mathcal P_K\mathcal L_{s_0}\phi\|_\infty.\] The printed conclusion is \[
 s_0+\frac{\varepsilon^--\mathrm{err}}{D^-\phi^+}
 <s_*<
 s_0+\frac{\varepsilon^++\mathrm{err}}{D^+\phi^-},
 \qquad \rho(\mathcal L_{s_*})=\lambda.\] The proof treats a forward parameter displacement and then describes the lower bound as analogous. A negative displacement reverses the relevant integral inequality. Dividing a negative residual by the larger bound \(D^-\phi^+\) can move the proposed lower endpoint to the wrong side of the root.\par
\subsection*{Counterexample satisfying the assumptions}
Fix \(d=1\), any Bernstein ellipse radii \(0<r<R\), and \(p=2\). Use a single analytic branch and weight \[v(z)=0,\qquad w_s(z)=2-s,\qquad
 s\in[1/2,3/2],\qquad
 (\mathcal L_sf)(z)=(2-s)f(0).\] The branch maps the larger ellipse into the smaller one. The weight is holomorphic and bounded, and the family is differentiable in \(s\). The operator is positive and rank one, with spectrum \(\{2-s,0\}\), a simple positive leading eigenvalue and a spectral gap for every parameter in the specified interval. Thus the analytic, summability, positivity and spectral assumptions all hold.\par
For every strictly positive continuous test function \(\psi\), \[\frac{d}{ds}\mathcal L_s\psi=-\psi(0),\qquad
 -2\sup\psi\le-\psi(0)\le-\inf\psi<0.\] Consequently [derivative] (\#derivative) holds with \(D^-=2\), \(D^+=1\). These are valid bounds; the theorem imposes no optimality condition on them.\par
Now choose \[\phi\equiv1,\quad \phi^-=\phi^+=1,\quad K=1,\quad
 \lambda=1,\quad s_0=\frac54,\quad
 \varepsilon^-=-\frac13,\quad \varepsilon^+=\frac18.\] Chebyshev interpolation preserves constants, so \(\mathrm{err}=0\). The actual residual is \[\mathcal P_1\mathcal L_{5/4}\phi-\phi=-\frac14,
 \qquad -\frac13<-\frac14<\frac18.\] The unique exact root is \(s_*=1\), because \(\rho(\mathcal L_s)=2-s\). In contrast, [printed] (\#printed) asserts \[\frac{13}{12}
 =\frac54+\frac{-1/3}{2}
 <s_*<
 \frac54+\frac18=\frac{11}{8}.\] Both proposed endpoints are within the original parameter interval, but the asserted lower endpoint exceeds the true root. This is a counterexample to the theorem, not a floating-point discrepancy. A separately authored 17-line Python check using exact `fractions.Fraction` arithmetic reproduces these fractions; the operator and domain assumptions are verified analytically above.\par
\subsection*{Static implementation evidence and counterevidence}
The rigorous Julia source sets \(D^-=33/10\) and \(D^+=59/100\) at lines 27–28, as in Proposition 6.5. Its function `dimensionbound`, lines 340–350, returns at line 349:\par
> s + hull(minmaxest.lo/(Dminus*phi\_plus),
> minmaxest.hi/(Dplus*phi\_minus))\par
where the displayed ASCII names stand for the corresponding Greek-named variables in the source. The lower residual is divided by \(D^-\phi^+\) regardless of its sign. Taking the hull of the two selected quotients cannot restore a missing, smaller quotient. For a valid residual enclosure \([-1/3,1/8]\) and denominator interval \([1,2]\), this rule again produces \([13/12,11/8]\) after adding \(5/4\).\par
This static match is evidence that the error is not confined to an isolated printed sign. It does not establish the sign, width or value of the actual residual in the publication-time run. The README’s stated output interval is a repository assertion, not an execution receipt from this review. The printed Algorithm 3, p. 944 (PDF p. 36), also contains subtraction signs that differ from the theorem and source return expression; it was not used to infer the execution history.\par
\subsection*{A precise repair}
Equation reference: repair. Let \[\ell=\varepsilon^--\mathrm{err},\quad u=\varepsilon^++\mathrm{err},\qquad
 m=D^+\phi^->0,\quad M=D^-\phi^+\ge m.\] Assume that the root exists in the stated parameter interval, as in the original theorem, and retain its positive-operator min–max implications. A valid enclosure is \[
 s_0+\min\{\ell/m,\ell/M\}
 <s_*<
 s_0+\max\{u/m,u/M\}.\] Equivalently, use the full interval quotient \([\ell,u]/[m,M]\). Endpoints outside the known parameter interval can be intersected with that interval. The theorem’s original hypotheses alone should not be read as an additional existence proof for a root at an arbitrary \(\lambda\).\par
Verification of the repair. Put \(F_s(x)=\mathcal L_s\phi(x)-\lambda\phi(x)\). Then \(-M\le\partial_sF_s\le-m\) and \(\ell<F_{s_0}(x)<u\). For a nonnegative displacement \(\delta\), \[F_{s_0}(x)-M\delta\le F_{s_0+\delta}(x)
 \le F_{s_0}(x)-m\delta.\] For a negative displacement, \[F_{s_0}(x)-m\delta\le F_{s_0+\delta}(x)
 \le F_{s_0}(x)-M\delta.\] Thus at \(\delta_-=\ell/M\) when \(\ell\ge0\), or \(\delta_-=\ell/m\) when \(\ell<0\), one has \(F_{s_0+\delta_-}>0\). At \(\delta_+=u/m\) when \(u\ge0\), or \(\delta_+=u/M\) when \(u<0\), one has \(F_{s_0+\delta_+}<0\). The min–max comparison and monotonicity place the root between these parameters whenever the comparison endpoints lie in the parameter interval; an endpoint beyond its outer edge gives a trivial bound from the known location of the root. These are exactly the endpoints in [repair] (\#repair). ◻\par
For the counterexample the repaired interval is \([11/12,11/8]\), which contains \(1\). In the common case \(\ell<0<u\), both safe endpoints use the smaller denominator \(D^+\phi^-\). This is a repair of the enclosure theorem. It does not by itself revalidate the published numerical interval: that requires the actual residual, eigenfunction bounds and error quantities, followed by outward-rounded evaluation of the corrected endpoints.\par
\subsection*{Core impact}
Theorem 4.2 is expressly presented as a general contribution in the introduction. Section 7 then invokes it to transform validated transfer operator errors into the interval asserted in Theorem 1.1. The counterexample invalidates that stated general guarantee, and the same rule appears in the acquired rigorous source. This warrants Major under the specified calibration even if a corrected run preserves all reported digits. No such run was performed here. The mathematical value of the Apollonian dimension and the numerical truth of the particular printed interval have not been disproved.\par
\section*{Problem 2: analytic functions are not dense in the stated Hölder norm}
Judgment: Major core proof gap; analytic repair verified. The final paragraph of the proof of Proposition 2.3, printed p. 917 (PDF p. 9), takes \[B_0=B_1=\mathcal H^{2,2}_{R_A},\qquad B_2=C^\alpha([-1,1]^2)\] and asserts density in \(B_2\) before applying Baladi–Tsujii, Lemma A.1 [BT]. That lemma requires density in both Banach norms (author-manuscript PDF pp. 36–37). Density in the uniform norm is insufficient.\par
\subsection*{Complete obstruction under the usual convention}
For \(0<\alpha<1\), let \(X=[-1,1]^2\) and use the usual norm \[\|f\|_{C^\alpha}=\|f\|_\infty+
 \sup_{x\ne y}\frac{|f(x)-f(y)|}{\|x-y\|^\alpha}.\] The function \(f(x_1,x_2)=|x_1|^\alpha\) belongs to this space since \(\big||a|^\alpha-|b|^\alpha\big|\le|a-b|^\alpha\). Any \(g\in\mathcal H^{2,2}_{R_A}\) is differentiable near the origin, so \(g(t,0)-g(0,0)=O(t)\). Consequently \[\lim_{t\downarrow0}
 \frac{|(f-g)(t,0)-(f-g)(0,0)|}{t^\alpha}=1.\] It follows that \(\|f-g\|_{C^\alpha}\ge1\) for every such \(g\). The asserted density is therefore false.\par
An explicit little-Hölder convention would change the space, but the paper calls \(C^\alpha\) the space of Hölder continuous functions and does not state that convention. It cannot be silently supplied to validate this invocation. The following repair avoids that ambiguity.\par
\subsection*{Replacement proof with the original inputs}
Retain the original proof’s Hölder Ruelle–Perron–Frobenius conclusion: \(T=\mathcal A_s\) has a positive leading eigenvalue \(\rho>0\), a nonnegative nonzero eigenfunction \(h\), and a spectral gap. Normalize \(\|h\|_\infty=1\). Also retain compactness of the operator \(L=\mathcal A_s\) on the Hardy space, which the original argument establishes. The constant function \(1\) lies in that space with norm one. Positivity gives \[T^n1\ge T^nh=\rho^nh,\qquad
 \|L^n\|\ge\|L^n1\|_{\mathcal H}\ge\|T^n1\|_\infty\ge\rho^n.\] The spectral-radius formula yields \(r_{\mathcal H}(L)\ge\rho\).\par
If \(r_{\mathcal H}(L)>\rho\), compactness, after complexification if necessary, supplies a nonzero analytic eigenfunction with eigenvalue \(\mu\) satisfying \(|\mu|=r_{\mathcal H}(L)>\rho\). Restriction to the real square is nonzero by analytic uniqueness and belongs to \(C^\alpha(X)\); it is therefore a Hölder eigenfunction with eigenvalue \(\mu\), contradicting the Hölder spectral-radius bound. Hence \(r_{\mathcal H}(L)=\rho\). Compactness and the Hölder spectral gap identify the Hardy leading eigenvalue as the same positive eigenvalue. Restriction also injects its generalized eigenspace into the Hölder generalized eigenspace, preserving simplicity.\par
This restores the required spectral identification and therefore the original monotonicity and dimension correspondence, conditional on the same RPF and compactness inputs already used by the authors. It does not reprove the underlying graph-directed-system theory or the computer-assisted mapping bounds. The supported error is the printed density argument; Proposition 2.3 itself has not been refuted. Because that proposition connects the computed Hardy spectrum to the dimension, the gap is in a core dependency rather than a peripheral citation. Its repair is already analytically verified here.\par
\section*{Checked reasoning, alternative readings and remaining work}
The finite-rank strategy is coherent: summable analytic weights give the operator norm bound, interpolation controls the difference on the real cube, and a positive trial eigenfunction allows pointwise residual bounds to constrain the spectral radius. The two-dimensional \(p=2\) coefficient-decay argument yields \(e^{-R\|\mathbf k\|_2}\), and the shell-counting integral gives the stated \(E_{2,K}\) expression. This review did not certify all endpoint cases of the separate general-\(d,p\) formulation.\par
An apparent Euler–Maclaurin sign problem was checked against counterevidence and was not retained as an additional substantive finding. The code uses the standard negative Bernoulli derivative correction (`apollonian-rigorous.jl`, line 273; `apollonian-nonrigorous.jl`, line 181). The opposite-sign Bernoulli–zeta identity printed on p. 934 provides further context against reading one isolated sign as proof of an incorrect executed tail calculation. Likewise, unavailable optional material and non-rigorous preliminary eigenvectors are not defects: the intended second stage is a validated calculation, and the present finding addresses its final enclosure rule.\par
The public arXiv record checked on 25 October 2026 lists only v1. The inspected publisher record and targeted correction search did not identify a linked erratum resolving these concerns; this is the scope of that check, not proof that no correction exists anywhere. The acquired public repository revision is later than online publication and is kept separate from the publication-time execution provenance.\par
\section*{Conclusion and disposition}
Item | Judgment | Repair and remaining work\par
Theorem 4.2 and source return rule | Major; theorem falsified | Signed interval-division repair proved; actual high-precision certificate not recomputed.\par
Proposition 2.3 density step | Major proof gap under usual Hölder convention | Replacement proof verified from original RPF and compactness inputs.\par
Published 128-place interval | Numerical validity unresolved here | No counterexample to the actual dimension; original residual data and corrected outward-rounded certification are needed.\par
Adoption and acceptance | Human decision pending | No author contact or adoption; focused independent checks completed; human review record empty.\par
The initial-review outcome is Major, supported by explicit mathematics and fixed-version source evidence. The demonstrated consequence is a defect in the claimed certificate and a repairable spectral proof gap. A corrected numerical certificate may retain the reported precision; the present work supplies no result about that recomputation.\par
9 P. L. Vytnova and C. L. Wormell, Hausdorff dimension of the Apollonian gasket, Invent. Math. 239 (2025), 909–946. Publisher version of record. https://doi.org/10.1007/s00222-024-01311-y. C. L. Wormell, apollonian, public repository, commit `4f73a6a1807ea7cb28313c379ffeeaf1410ff1d6`, 13 February 2025. README and two Julia files, static inspection. https://github.com/wormell/apollonian. V. Baladi and M. Tsujii, Dynamical determinants and spectrum for hyperbolic diffeomorphisms, Contemp. Math. 469 (2008), 29–68. Consulted author manuscript: arXiv:math/0606434v2, 17 January 2008, Appendix A, Lemma A.1, PDF pp. 36–37. https://arxiv.org/abs/math/0606434v2. P. Vytnova and C. Wormell, Hausdorff dimension of the Apollonian gasket, arXiv:2406.04922v1, 7 June 2024. Version-history check only; not substituted for the journal PDF. https://arxiv.org/abs/2406.04922.\par
\section*{Supporting checks}
The publisher PDF SHA-256 is\par
c6c5aecaed0266668aa654cacc8b86827882679f2234974d0d18659f75459b58.\par
The source-manifest SHA-256 is\par
1861b291e4e8f497f8f278670bd7756819d8ddd705339ba65a6f6037218a184d.\par
Its output is \[\mathrm{residual}=-1/4,\quad
 \mathrm{printed}=[13/12,11/8],\quad
 \mathrm{repaired}=[11/12,11/8],\quad
 s_*=1.\] This computation checks fractions; it is not a substitute for the analytic proof of assumptions above.\par
Reusable publisher page images are retained for PDF pages 2, 3, 9, 12, 16, 18, 20, 21, 26, 33, 35, 36 and 37. The reference’s PDF pages 36–37 are retained separately. Text extraction, rendering commands, image hashes and inspection notes are recorded in the evidence and provenance files. The source PDF produced a Poppler structural warning while parsing; extraction and rendering succeeded. No mathematical conclusion is based on that warning.\par
They were not independent acceptance reviewers. The session’s separate runtime self-description is GPT-6 Astra Pro; independent backend verification is unavailable and no equivalence between these labels is asserted. No model change was made.\par
The exact TeX source, compiler invocation, compiler executable hash, dependency/cache input hashes, logs, final PDF hash and report-page inspection record accompany this report in the assigned staging tree. Compilation uses the already authorized Tectonic 0.17.0 with `–untrusted`; shell escape and author-code execution are disabled. The reported findings have undergone focused independent acceptance. Human reviews remain empty and scientific adjudication remains pending.\par
\end{document}
